\newpage
\phantomsection
\label{prf:two-sided-induction-and-recursion}
\LRAProofFor{thm:two-sided-induction-and-recursion}

\begin{remark*}[Return]
\hyperref[thm:two-sided-induction-and-recursion]{Return to Theorem}
\end{remark*}

\begin{theorem*}[Two-sided induction and recursion]
A predicate that holds at zero and is preserved by successor and predecessor
holds everywhere on \(Z\). Compatible forward and backward recursive clauses
determine a unique function out of \(Z\).
\end{theorem*}

\begin{proof}[Professional Standard Proof]
\LRAProofBodyStart
For induction, use the disjoint decomposition
\(Z=\{0\}\sqcup P\sqcup N\). The zero case is the base hypothesis. A point of
\(P\) is obtained by starting at \(\mathsf{succ}(0)\) and repeatedly
applying successor, so ordinary induction along \(P\) proves the predicate on
the positive ray. A point of \(N\) is obtained by starting at
\(\mathsf{pred}(0)\) and repeatedly applying predecessor, so ordinary
induction along \(N\) proves the predicate on the negative ray.

For recursion, define the function separately on the three parts: send zero to
the chosen base value, define the positive ray by iterating the successor
step, and define the negative ray by iterating the predecessor step. The
inverse hypotheses on the two step functions guarantee that the successor and
predecessor equations agree at the zero crossings. Uniqueness follows by
two-sided induction applied to the predicate that two candidate functions have
the same value.
\end{proof}

\begin{proof}[Detailed Learning Proof]
\LRAProofBodyStart
This is the two-sided analogue of Peano induction. Peano induction has one
direction away from the base point. The integer line has two directions away
from zero, so a proof must propagate both right and left. The recursion theorem
has the same shape: defining a function on integers requires a zero value, a
right-step rule, and a left-step rule. The two step rules must be compatible,
because the first positive and first negative points both touch zero through
the inverse laws.
\end{proof}

\begin{remark*}[Proof structure]
Use ordinary structural induction on the positive and negative rays; prove
uniqueness by applying the resulting two-sided induction principle to equality
of two recursive candidates.
\end{remark*}

\begin{dependencies}
\begin{itemize}
  \item \hyperref[thm:two-sided-successor-predecessor-inverses]{Successor and predecessor are inverse operations}
\end{itemize}
\end{dependencies}

\clearpage
